\section{导读：这期为什么是多模态技术路线课}

本节先建立阅读框架。这期是张小珺商业访谈录里少见的高技术密度访谈，李广密作为 co-host，把问题集中在多模态十年史、视觉与语言的差异、o1 范式对多模态的启发，以及未来两年可能出现的 GPT-4 时刻。张祥雨的叙述不是一组论文八卦，而是一条技术路线：CV 曾经持续向 NLP 学习，但静态图像本身和自然语言不同；多模态真正的突破，可能要从 next token prediction、RL、CoT、动作空间和在线学习重新理解。

\begin{importantbox}{本期核心命题}
多模态的难点不只是“把图像、视频和文本放进一个模型”。静态图像里生成、理解和人类对齐天然分离；视觉生成与视觉理解长期没有真正融合；o1 之后，研究者开始意识到：多模态也可能需要思维链、反思、动作空间扩展和环境反馈，才能出现真正的 GPT-4 时刻。
\end{importantbox}

\begin{knowledgebox}{视觉策略说明}
本视频是固定访谈画面，没有 slides 或屏幕演示。正文只使用封面作为来源识别；正文图像全部为自制概念图，用来解释 CV/NLP 迁移史、静态图像三重割裂、生成理解 gap、next token 缺陷、RL/CoT、视觉 Long CoT、long context 多模型协作和在线学习。
\end{knowledgebox}

\begin{figure}[H]
\centering
\includegraphics[width=0.96\textwidth]{cv-to-nlp-learning-history.png}
\caption{CV 向 NLP 学习的历史：从 ImageNet/ResNet 到 ViT、iGPT、MIM 与多模态大模型。自制概念图，依据 00:02:00--00:17:14 对谈内容整理。}
\end{figure}

\begin{knowledgebox}{读图：CV 的 GPT 时刻没有简单复制 NLP}
图中从 ImageNet、ResNet 到 ViT/iGPT/MIM，再到 VLM，展示的是 CV 不断借鉴 NLP 范式的过程。关键问题是：迁移架构容易，迁移数据属性和目标函数很难。语言数据天然来自人类表达，静态图像却不天然包含人类理解。
\end{knowledgebox}

\subsection{阅读路线}

为了避免被术语淹没，本节给出读这期的四个问题。第一，为什么 CV 向 NLP 学了十年，却没有在静态图像上复刻 GPT？第二，为什么图像生成和图像理解放进一个系统后仍然像两个模型？第三，为什么大模型越大，文科和对话能力更强，但数学推理可能会跳步、退化？第四，为什么 o1 的反思和 CoT pattern 会让多模态研究重新看到路？

\noindent\begin{tabularx}{\textwidth}{p{0.19\textwidth}p{0.34\textwidth}X}
\toprule
阅读问题 & 访谈中的材料 & 要形成的判断 \\
\midrule
CV 为什么难复制 GPT & 静态图像、生成/理解/对齐割裂 & 数据形态决定学习目标，不能只迁移架构。 \\
多模态为何未融合 & 理解模型和生成模型各自变强但互不增益 & 缺少共同的思考过程和可控动作空间。 \\
推理为何会退化 & 大模型数学跳步、next token 压缩缺陷 & 分布压缩不等于计算精度。 \\
未来突破在哪里 & RL、CoT、视觉 Long CoT、在线学习 & 关键在动作空间、反馈和持续学习。 \\
\bottomrule
\end{tabularx}

\subsection{本章小结}

EP102 的主线是：多模态不是“多加一种输入”这么简单。要理解未来两年的 GPT-4 时刻，必须同时理解静态图像数据的特殊性、语言模型目标函数的缺陷、o1 引入的反思模式，以及在线学习对下一代 Agent 的意义。

